% ==============================================================================
% CAPÍTULO 3 - REFERENCIAL TEÓRICO
% ==============================================================================
    
\section{Referencial Teórico} 
\label{sec:referencial-teorico}

Esta seção organiza o estado da arte em cinco eixos que sustentam o desenho do
experimento: a avaliação de agentes de código em repositórios reais, a
engenharia de contexto e os arquivos de instrução de repositório, as
ferramentas que executam esses agentes, o uso de LLMs como avaliadores e as
métricas de qualidade de código. Ao fim de cada eixo, indica-se o que este
trabalho acrescenta.

\subsection{Benchmarks de agentes de código em repositórios reais}
\label{sec:ref-benchmarks}

A avaliação de geração de código por LLMs começou com problemas isolados.
\citeonline{chen2021codex} introduziram o HumanEval e a métrica
\textit{pass@k}, a probabilidade de que ao menos uma entre $k$ amostras passe
nos testes do problema, e \citeonline{austin2021mbpp} seguiram o mesmo formato
com o MBPP. \citeonline{liu2023evalplus} mostraram, porém, que os testes
desses conjuntos eram insuficientes: ao ampliá-los no EvalPlus, a taxa de acerto
de diversos modelos caiu de forma expressiva, o que evidencia a dependência de
qualquer \textit{benchmark} da qualidade do seu oráculo.

O passo para tarefas em escala de repositório foi dado pelo SWE-bench
\cite{jimenez2024swebench}, que reúne \textit{issues} reais de projetos Python
no GitHub e avalia a correção pelos testes do PR que as resolveu, separando
testes que devem passar a falhar-para-passar dos que devem continuar passando.
A versão SWE-bench Verified \cite{openai2024swebenchverified} filtrou
manualmente os itens com enunciados subespecificados ou testes injustos,
reconhecendo que parte das falhas medidas vinha do \textit{benchmark} e não dos
modelos. Sobre essa base, \citeonline{yang2024sweagent} mostraram, com o
SWE-agent, que a interface entre o agente e o computador (comandos de busca,
edição com verificação sintática, visualização paginada de arquivos) altera
substancialmente a taxa de resolução para um mesmo modelo, e
\citeonline{wang2024openhands} consolidaram no OpenHands uma plataforma aberta
para agentes generalistas de desenvolvimento. Em paralelo, trabalhos de
completação em nível de repositório, como o de \citeonline{shrivastava2023repo}
e o RepoCoder de \citeonline{zhang2023repocoder}, demonstraram que selecionar
contexto relevante do restante do repositório melhora a geração, o que
fundamenta a ideia de que o repositório pode e deve informar o modelo.

Esses trabalhos têm em comum o uso de repositórios públicos, muitas vezes
presentes nos dados de treinamento dos modelos avaliados, e a concentração em
Python. Este trabalho acrescenta um \textit{benchmark} construído sobre um
repositório privado em PHP/Laravel, cujos \textit{commits} não são acessíveis
aos modelos, com o mesmo princípio de oráculo oculto do SWE-bench, e usa esse
\textit{benchmark} não para ordenar modelos, mas para medir, por ablação, o
efeito de uma intervenção no repositório: o \textit{harness}.

\subsection{Engenharia de contexto e arquivos de instrução de repositório}
\label{sec:ref-contexto}

Agentes de código seguem, em essência, o laço de raciocínio e ação proposto por
\citeonline{yao2023react}: o modelo alterna entre pensar, chamar uma ferramenta
e observar o resultado, acumulando tudo na janela de contexto. Como essa janela
é finita e o desempenho degrada quando a informação relevante fica dispersa ou
distante \cite{liu2024lost}, a seleção do que entra no contexto passou a ser
uma disciplina própria, chamada de engenharia de contexto
\cite{anthropic2025context}. A prática recomendada é separar o que deve estar
sempre presente (regras curtas, comandos de construção e teste, convenções) do
que deve ser lido sob demanda, quando a tarefa o exigir.

Essa separação materializou-se em arquivos versionados no repositório. O
\texttt{CLAUDE.md} do Claude Code é lido automaticamente no início de cada
sessão e deve conter o mínimo necessário para o agente operar no projeto
\cite{anthropic2025bestpractices}; o \texttt{AGENTS.md} propõe-se como formato
aberto e neutro em relação à ferramenta para o mesmo fim \cite{agentsmd2025}.
\textit{Skills} são documentos mais longos, organizados por assunto, que o
agente abre quando a tarefa toca aquele assunto; a sua existência só é anunciada
nas regras sempre ativas. A configuração de ferramentas, por fim, habilita
servidores de linguagem e outros recursos que mudam o que o agente consegue
observar sem ler arquivos inteiros.

A literatura sobre esses artefatos é, em grande parte, prescritiva: descreve o
que colocar neles, mas raramente mede o efeito em condições controladas
\pendente{CITAÇÃO: estudos empíricos sobre arquivos de instrução de repositório
e sua eficácia}. Este trabalho acrescenta uma ablação em quatro níveis
cumulativos, aplicada a um \textit{harness} real e congelado, com o mesmo
enunciado e o mesmo ambiente em todos os níveis, de modo a atribuir qualquer
diferença observada à camada adicionada.

\subsection{Ferramentas e agentes de código}
\label{sec:ref-ferramentas}

O Claude Code \cite{anthropic2025claudecode} e o OpenCode
\cite{opencode2025docs} são agentes de terminal: recebem uma tarefa em
linguagem natural, dispõem de ferramentas de leitura, busca, edição e execução
de comandos, e trabalham até concluir. Ambos leem arquivos de instrução do
repositório e podem ser executados em modo não interativo, o que os torna
adequados a um experimento automatizado. Eles diferem no \textit{prompt} de
sistema, no conjunto de ferramentas e na gestão de contexto, o que torna a
comparação entre modelos executados por \textit{runners} distintos uma
comparação de pares modelo-\textit{runner}, e não de modelos isolados.

A escolha do modelo, por sua vez, envolve uma relação entre capacidade e
custo descrita pelas leis de escala: o desempenho melhora de forma previsível
com parâmetros, dados e computação \cite{kaplan2020scaling}, e há um tamanho
ótimo de modelo para cada orçamento de treinamento \cite{hoffmann2022chinchilla}.
Para quem usa os modelos, isso se traduz em famílias com variantes de tamanhos e
preços distintos e em modelos abertos, como o gpt-oss-20b
\cite{openai2025gptoss}, que podem ser executados localmente com
\textit{runtimes} como o llama.cpp \cite{llamacpp2025}, com quantização e as
limitações de \textit{hardware} que isso implica.

Este trabalho acrescenta uma comparação do mesmo \textit{harness} em seis
modelos, de uma família comercial de ponta a um modelo aberto quantizado em
GPU de consumo, executados por dois \textit{runners} sob um mesmo protocolo,
com a comparação de níveis sempre feita dentro de cada par modelo-\textit{runner}.

\subsection{LLM como avaliador}
\label{sec:ref-juiz}

Quando a métrica de interesse não é binária, como a qualidade de um
\textit{pull request}, a avaliação humana é cara e a avaliação automática
tradicional é limitada. \citeonline{zheng2023judging} propuseram usar LLMs
fortes como juízes e mostraram, com o MT-Bench e a Chatbot Arena, que o GPT-4
concorda com avaliadores humanos em nível comparável ao da concordância entre
humanos. O mesmo trabalho catalogou vieses do método: viés de posição (a
preferência pela resposta apresentada primeiro), viés de verbosidade (a
preferência por respostas mais longas) e autopreferência (a preferência por
respostas do próprio modelo). \citeonline{wang2024fair} mostraram que a simples
troca de ordem das respostas pode inverter o veredito, e
\citeonline{panickssery2024self} demonstraram que a autopreferência está ligada
à capacidade do modelo de reconhecer os próprios textos. Para o viés de
verbosidade, \citeonline{dubois2024length} propuseram controlar o comprimento
na análise. No domínio de código, \citeonline{zhuo2024icescore} mostraram que
LLMs instruídos com critérios explícitos avaliam utilidade e correção funcional
com correlação superior à de métricas baseadas em similaridade textual.

Este trabalho aplica essas lições ao desenho dos juízes: usa dois juízes por
mudança, nunca da mesma família do modelo que escreveu o código, para mitigar a
autopreferência; avalia cada mudança isoladamente, com nota absoluta e
evidência obrigatória por dimensão, para evitar o viés de posição; instrui
explicitamente a não recompensar o comprimento; e, sobretudo, mede a
concordância entre juízes e a correlação das notas com medidas independentes
(testes ocultos, regras de convenção e análise estática) em vez de assumir a
sua validade. A ligação entre um juiz (Muse Spark) e o repositório-sujeito é
discutida como ameaça à validade.

\subsection{Métricas de qualidade de código}
\label{sec:ref-qualidade}

Qualidade de código é um construto multidimensional. A norma ISO/IEC~25010
\cite{iso25010} organiza a qualidade de produto em características como
adequação funcional, segurança e manutenibilidade, e esta última desdobra-se
em analisabilidade, modificabilidade e testabilidade. Medidas estruturais
tentam aproximar essas características: a complexidade ciclomática de
\citeonline{mccabe1976} conta os caminhos linearmente independentes de uma
função e é usada há décadas como indicador de dificuldade de teste e de
manutenção. A análise estática complementa essas medidas ao apontar erros de
tipo, chamadas a métodos inexistentes e outros defeitos sem executar o código;
no ecossistema PHP, o PHPStan e sua extensão para Laravel, o Larastan
\cite{larastan2025}, cumprem esse papel em níveis de rigor crescentes.

Para código gerado por LLMs, a literatura tem apontado que correção funcional
não implica qualidade. \citeonline{pearce2022asleep} encontraram
vulnerabilidades em cerca de 40\% dos cenários em que o GitHub Copilot produziu
código relevante para segurança, o que motiva uma dimensão de segurança
separada da de requisitos. Práticas de engenharia como o desenvolvimento
orientado a testes \cite{beck2002tdd} e os princípios de legibilidade e
simplicidade de \citeonline{martin2008clean} fornecem o vocabulário das
convenções que os repositórios codificam em seus arquivos de instrução e que
revisores humanos cobram em \textit{pull requests}.

Este trabalho acrescenta uma operacionalização em três camadas que combina
medidas objetivas (Larastan, complexidade, testes escritos), regras
heurísticas derivadas das convenções do próprio repositório e julgamento por
LLM em sete dimensões alinhadas à ISO/IEC~25010 e às convenções do projeto,
com verificação cruzada entre as camadas. A escolha de medir o que o
\textit{harness} prescreve (testes, convenções, autorização, simplicidade) é
deliberada: é nessas dimensões que uma intervenção textual tem como agir.
